\documentclass[11pt,a4paper]{article}
\usepackage[margin=26mm]{geometry}
\usepackage[T1]{fontenc}
\usepackage{lmodern,microtype}
\usepackage{amsmath,amssymb,amsthm,mathtools}
\usepackage{xcolor}
\usepackage[colorlinks=true,linkcolor=blue!45!black,citecolor=blue!45!black,urlcolor=blue!45!black]{hyperref}
\usepackage{fancyhdr}
\setlength{\headheight}{14pt}
\pagestyle{fancy}
\fancyhf{}
\fancyhead[L]{\small A geometric argument for an arithmetic quadric}
\fancyhead[R]{\small Research draft}
\fancyfoot[C]{\thepage}
\setlength{\parskip}{2pt}
\setlength{\parindent}{14pt}
\newtheorem{theorem}{Theorem}[section]
\newtheorem{lemma}[theorem]{Lemma}
\newtheorem{proposition}[theorem]{Proposition}
\theoremstyle{remark}
\newtheorem{remark}[theorem]{Remark}
\newcommand{\Z}{\mathbf Z}
\newcommand{\Q}{\mathbf Q}
\newcommand{\R}{\mathbf R}
\newcommand{\C}{\mathbf C}
\newcommand{\F}{\mathbf F}
\newcommand{\PP}{\mathbf P}
\newcommand{\OO}{\mathcal O}
\newcommand{\m}{\mathfrak m}
\newcommand{\n}{\mathfrak n}
\newcommand{\Spec}{\operatorname{Spec}}
\newcommand{\Frac}{\operatorname{Frac}}
\newcommand{\Pic}{\operatorname{Pic}}
\newcommand{\Cl}{\operatorname{Cl}}
\newcommand{\divisor}{\operatorname{div}}
\newcommand{\Gal}{\operatorname{Gal}}
\newcommand{\colim}{\operatorname*{colim}}
\newcommand{\HD}{H_{\mathcal D,\mathrm{an}}}
\newcommand{\reg}{\operatorname{reg}}
\begin{document}
\thispagestyle{plain}
\begin{center}
{\Large An arithmetic quadric and rational Gersten injectivity}\\[4pt]
{\normalsize A proposed counterexample by a geometric regulator argument}\\[8pt]
{\small Research draft --- 7 October 2026}
\end{center}
\vspace{3pt}
\noindent\fcolorbox{black!35}{black!3}{%
  \parbox{\dimexpr\textwidth-2\fboxsep-2\fboxrule\relax}{\small
  \textbf{Review status.} This note gives a new proposed counterexample,
  with the argument and its standard inputs specified below. It has
  undergone internal adversarial checking, but has not been externally
  refereed or formally verified. The cited results establish the tools;
  they do not establish the new conclusion claimed here.}}

\begin{abstract}
For an odd prime \(p\), consider the regular local ring
\(\bigl(\Z_{(p)}[x,y,z]/(p+xy+z^2)\bigr)_{(x,y,z)}\).
A Borel coefficient on its regular divisor \(x=0\) gives an actual
supported \(K_3\)-class. After a ramified quadratic base change, this
class becomes a coefficient times a ruling of a split affine quadric.
A small resolution proves that the ruling persists arithmetically.
A componentwise Galois argument then shows that its degree-two Betti
class survives every permitted finite localization. The multiplicative
analytic Deligne regulator detects the resulting class at every such
stage. The proposed conclusion is an infinite-order integral class
with zero generic image.
\end{abstract}

\section{The ring and the supported coefficient}

We use Quillen algebraic \(K\)-theory; write \(M_{\Q}=M\otimes_{\Z}\Q\).
Pushforwards from regular closed subschemes may be defined in
\(G\)-theory and transported to \(K\)-theory on the regular schemes
involved. We use localization, d\'evissage, flat base change and the
projection formula \cite{Quillen}.

\begin{theorem}[Proposed counterexample]\label{thm:geometricmain}
For every odd prime \(p\), the ring
\[
 A=\left(\Z_{(p)}[x,y,z]/(p+xy+z^2)\right)_{(x,y,z)}
\]
is a regular local domain of dimension three and residue field \(\F_p\).
There is an integral class \(c\in K_3(A)\) of infinite order whose
image in \(K_3(\Frac A)\) is zero. Consequently
\[
 \ker\!\left(K_3(A)_{\Q}\longrightarrow K_3(\Frac A)_{\Q}\right)\ne0.
\]
\end{theorem}

Put
\[
 V=\Z_{(p)},\qquad W=V[u]/(u^2+p),\qquad E=\Q(u),\qquad
 \sigma(u)=-u.
\]
The ring \(W\) is a DVR with uniformizer \(u\) and residue field
\(\F_p\); \(E\) is imaginary quadratic. Fix an embedding
\(E\hookrightarrow\C\). We use the analytic Deligne regulator, with
\(\R(j)=(2\pi i)^j\R\). Its point component is
\[
 r_2:K_3(E)\longrightarrow
 \HD^1(\mathrm{pt},\R(2))=\C/\R(2).
\]
Borel's theorem and its comparison with the Beilinson regulator imply
that \(r_2\) is nonzero \cite[Proposition~12.2]{Borel};
\cite[Section~6, especially (6.4)]{BunkeTamme}.
The precise regulator scope used on varieties is explained in
Section~\ref{sec:geometricregulator}.

The ambient local ring \(V[x,y,z]_{(p,x,y,z)}\) is regular of
dimension four. The equation \(p+xy+z^2\) has nonzero linear term
\(p\), so its quotient \(A\) is regular of dimension three.
It is a domain, with \(\Frac A=\Q(x,z)\), and
\(p=-xy-z^2\in\m^2\), where \(\m=(x,y,z)A\).
Its regular principal divisor has coordinate ring
\[
 R_D=A/(x)\cong W[y]_{(u,y)},\qquad u\longmapsto z.
\]
Write \(i:\Spec R_D\hookrightarrow\Spec A\).

Quillen's finite-field calculation gives \(K_2(\F_p)=0\)
\cite{QuillenFinite}; hence localization makes
\(K_3(W)\to K_3(E)\) surjective. Choose \(b_0\in K_3(W)\)
whose image has nonzero regulator, and set
\begin{equation}\label{eq:geometriccoefficient}
 \beta=b_0-\sigma b_0,\qquad
 c=i_*(\beta|_{R_D})\in K_3(A).
\end{equation}
Thus \(\sigma\beta=-\beta\) integrally. On \(\C/\R(2)\), complex
conjugation acts by \(-1\); therefore
\(r_2(\beta|_E)=2r_2(b_0|_E)\ne0\).
Localization at \(x\) sends \(c\) to zero, so its generic image is
already zero integrally. It remains to prove that \(c\) has infinite
order.

\section{A small resolution and arithmetic Picard persistence}

Write
\[
 A_0=V[x,y,z]/(p+xy+z^2),\qquad
 S=A_0\setminus(x,y,z),\qquad
 B_0=W[x,y,z]/(xy+z^2-u^2).
\]
Then \(B=A\otimes_VW=S^{-1}B_0\). Define
\[
 Q=\Spec E[x,y,z]/(xy+z^2-u^2),\qquad T=\Spec B[1/u].
\]
For \(s\in S\), let \(Q_s=D(s)\subset Q\).
Continuity for filtered localizations
\cite[Chapter~IV, Section~6.4]{Weibel} gives
\begin{equation}\label{eq:geometriccontinuity}
 K_3(T)=\colim_{s\in S}K_3(Q_s).
\end{equation}

The finite free \(A\)-algebra \(B\) is local: its closed fiber is
\(\F_p[u]/(u^2)\), with a unique maximal ideal. Consequently
\[
 B=(B_0)_{\n},\qquad \n=(u,x,y,z).
\]
It is a domain, since \(u^2+p\) stays irreducible over \(\Q(x,z)\).
Its generic fiber is smooth; off the origin of the special fiber
one of the partial derivatives \(y,x,2z\) is a unit.
Thus \(B\) is regular off its closed point and is normal by
the hypersurface \(S_2\) property and \(R_1\).
The height-one primes
\[
 P_+=(x,z-u),\qquad P_-=(x,z+u)
\]
occur with multiplicity one in \(\divisor_B(x)=P_++P_-\).

\begin{lemma}\label{lem:geometricresolution}
The class \([P_+]\) has infinite order in \(\Cl(B)\).
\end{lemma}
\begin{proof}
Set \(w=z-u\) and \(d=2u\), so the equation is
\(xy+w(w+d)=0\). Blow up \(I=(x,w)\) in \(\Spec B\).
The two actual Rees charts, inside the fraction field, are
localizations of
\begin{align*}
 W[x,t]:&\quad w=xt,\qquad y=-t(xt+d),\\
 W[q,y]:&\quad w=-qy-d,\qquad x=-q(qy+d).
\end{align*}
For example the first chart is \(B[w/x]\): substituting \(w=xt\)
gives \(x(y+t(xt+d))=0\), and \(x\)-saturation gives the displayed
polynomial chart. The other chart is obtained by \(w\)-saturation.
There are no further relations, since the indicated parameter pairs
are algebraically independent over \(E\).

The blowup \(Y\) is therefore regular. The inverse image of the
closed point has charts \(\F_p[t]\) and \(\F_p[q]\), glued by
\(q=t^{-1}\), so it is the reduced curve \(F\simeq\PP^1_{\F_p}\).
The denominators coming from \(B\) reduce to nonzero constants
on these charts, so no points of this projective line are removed.

The ideal \(I\) is invertible away from the closed point.
At a nonclosed prime containing \(I\), if \(u\) is not a unit
then \(y\) is a unit and \(I=(w)\). If \(u\) is a unit, then
\(w+d\) is a unit and \(I=(x)\). Outside \(V(I)\) it is the
unit ideal. Hence \(Y\to\Spec B\) is an isomorphism off the
closed point. Its exceptional locus is the curve \(F\), of
codimension two; there are no exceptional prime divisors.

The tautological identity \(I\OO_Y=\OO_Y(1)\)
\cite[Section~31.33]{StacksBlowup} restricts to
\(\OO_{\PP^1}(1)\) on \(F\).
The ideal \(I\) has valuation one at \(P_+\) and zero at all
other height-one primes. Since the resolution has no exceptional
divisors, for the strict transform \(\widetilde P_+\) we obtain
\[
 I\OO_Y=\OO_Y(-\widetilde P_+),\qquad
 \OO_Y(\widetilde P_+)|_F=\OO_{\PP^1}(-1).
\]
If \(\divisor_B(f)=N P_+\), the same rational function upstairs
has divisor exactly \(N\widetilde P_+\). Restricting its trivial
divisor line bundle to \(F\) forces \(\OO_{\PP^1}(-N)\) to be
trivial, and hence \(N=0\).
\end{proof}

The quotient
\[
 B/(u)=\F_p[x,y,z]_{(x,y,z)}/(xy+z^2)
\]
is a domain. Inverting \(u\) therefore removes only the principal
height-one prime \((u)\). The divisor-class localization sequence
gives an isomorphism
\begin{equation}\label{eq:geometricclasslocalization}
 \Cl(B)\xrightarrow{\;\sim\;}\Cl(B[1/u])=\Pic(T).
\end{equation}
In particular the image of \(P_+\) is not torsion.

\section{Every geometric boundary component has zero class}

Let \(D_\pm\subset Q\) be the divisors \((x,z\mp u)\).
The projective closure of \(Q\) is split:
\[
 \overline Q=\{xy+z^2=u^2t^2\}\cong\PP^1_E\times\PP^1_E,
 \qquad
 \det\begin{pmatrix}x&ut-z\\ut+z&y\end{pmatrix}=0.
\]
Its divisor at infinity is a smooth \((1,1)\) curve, and becomes
the diagonal after changing a factor coordinate. Thus
\begin{equation}\label{eq:geometricpic}
 Q\cong(\PP^1\times\PP^1)\setminus\Delta,\qquad
 \Pic(Q)=\Z^2/\Z(1,1)=\Z[D_+].
\end{equation}
The same calculation holds over \(\overline E\) and \(\C\).
The geometric generator is defined over \(E\), so the absolute
Galois group of \(E\) acts trivially on \(\Pic(Q_{\overline E})\).
Projection to a factor is an affine-line bundle; on complex points
it is a homotopy equivalence. In particular
\[
 H^2(Q_{\C},\Z)=\Z h,\qquad
 h=c_1^{\mathrm{top}}(\OO_Q(D_+))
\]
up to sign, and \(h\) is primitive.
By Lemma~\ref{lem:geometricresolution} and
\eqref{eq:geometricclasslocalization}, the generator in
\eqref{eq:geometricpic} has infinite-order image on \(T\).
Hence \(\Pic(Q)\to\Pic(T)\) is injective.

\begin{lemma}\label{lem:geometriccomponents}
For \(s\in S\), every geometric irreducible curve in
\(Q_{\C}\setminus(Q_s)_{\C}\) has zero class in \(\Pic(Q_{\C})\).
Consequently restriction is injective:
\begin{equation}\label{eq:geometricbetti}
 H^2(Q_{\C},\R)\longrightarrow H^2((Q_s)_{\C},\R).
\end{equation}
\end{lemma}
\begin{proof}
Let \(D\) be an \(E\)-prime divisor contained in \(V(s)\).
It is absent from \(T\), so its canonical Cartier section
trivializes \(\OO_Q(D)|_T\). Picard injectivity gives
\([D]=0\) on \(Q\).
In characteristic zero, \(D_{\overline E}\) is reduced and
its irreducible components form one Galois orbit.
Trivial Galois action on the geometric Picard group implies
that all these components have the same class \(a\in\Z\).
If there are \(r\) components, then \(ra=[D_{\overline E}]=0\);
thus \(a=0\). Base change to \(\C\) proves the first assertion.

For the second, let \(Z=V(s)_{\C,\mathrm{red}}\).
The kernel in \eqref{eq:geometricbetti} is the image of
\(H^2_Z(Q_{\C},\R)\).
Delete the finite set of singular and intersection points of
the curve components. Point-supported cohomology on a real
four-manifold vanishes below degree four, so neither
degree-two support cohomology nor the ambient \(H^2\) changes.
The remaining curves are disjoint and smooth.
The Thom isomorphism supplies one degree-two support generator
per component, mapping to that component's first Chern class.
Every such class is zero by the first assertion.
\end{proof}

\begin{remark}[An independent norm check]
A geometric boundary component \(\Gamma\) is defined over some finite
\(E'/E\). Set \(L=\OO_Q(D_+)\) and write
\([\OO(\Gamma)]=a[L_{E'}]\) in \(\Pic(Q_{E'})=\Z[L_{E'}]\).
This line bundle is trivial on \(T_{E'}\). Its finite-flat norm
(equivalently, determinants after restriction of scalars) therefore
makes \(L_T^{a[E':E]}\) trivial. Since \(L_T\) has infinite order,
\(a=0\). This also proves the componentwise assertion.
\end{remark}

The componentwise assertion is essential. Merely observing that
\(\divisor(s)\) is principal would not suffice:
\(\divisor(x)=D_++D_-\) is principal although its components have
opposite nonzero classes. Here \(x\notin S\).

\section{Regulator detection at every finite stage}
\label{sec:geometricregulator}

Throughout, \(\HD\) denotes \emph{analytic} Deligne cohomology.
Bunke--Tamme construct a natural multiplicative analytic Beilinson
regulator \cite[Remark~2.6 and Theorem~2.31]{BunkeTamme},
with weight-\(j\) component
\[
 r_j:K_a(U)\longrightarrow\HD^{2j-a}(U_{\C},\R(j)).
\]
Their site consists of regular separated schemes of finite type
over \(\Z\). To apply it here, spread any finite collection of
classes, opens and line bundles to a smooth model over
\(\OO_E[1/N]\) for a sufficiently divisible integer \(N\).
Such a model is in that site. Further integer localizations do not
change its complex fiber; naturality and \(K\)-theory continuity
give the stated maps on \(E\) and \(Q_s\). Projection to the chosen
embedding preserves multiplication. We use no regulator on \(T\).

\begin{lemma}\label{lem:geometricdeligne}
For every smooth complex surface \(U\),
\[
 \HD^3(U,\R(3))\cong H^2(U,\C)/H^2(U,\R(3)).
\]
These identifications are natural in \(U\), and the maps
\(\HD^3(Q_{\C},\R(3))\to\HD^3((Q_s)_{\C},\R(3))\)
are injective for every \(s\in S\).
\end{lemma}
\begin{proof}
The analytic Deligne complex is
\([\R(3)\to\OO_U\to\Omega_U^1\to\Omega_U^2]\).
Since there are no holomorphic three-forms, the holomorphic
Poincar\'e lemma identifies this with the shifted cone of
\(\R(3)\to\C\).
Real cohomology injects into its complexification, giving the
displayed quotient. An injection of real vector spaces
\(V\hookrightarrow W\) induces an injection
\[
 (V\otimes_{\R}\C)/(2\pi i)^3V
 \longrightarrow (W\otimes_{\R}\C)/(2\pi i)^3W.
\]
Apply Lemma~\ref{lem:geometriccomponents}.
\end{proof}

Put \(L=\OO_Q(D_+)\) and
\[
 \eta=[\OO_{D_+}]=1-[L^{-1}]\in K_0(Q),\qquad
 \gamma=(\beta|_E)\eta\in K_3(Q).
\]
Here \(\eta\) is the class of the two-term locally free resolution
of \(\OO_{D_+}\); it has rank zero and
\(\operatorname{ch}_1(\eta)=c_1(L)\).
Multiplicativity gives
\begin{equation}\label{eq:geometricproduct}
 r_3(\gamma)=r_2(\beta|_E)\cup c_1(L).
\end{equation}
There are no other weight-three terms: at a complex point,
positive-twist analytic Deligne cohomology is concentrated in
degree one, and the rank of \(\eta\) is zero.

The right side of \eqref{eq:geometricproduct} is nonzero.
Indeed the Betti normalization of \(c_1(L)\) is \(2\pi i\,h\).
Represent \(r_2(\beta|_E)\) by \(it\), \(t\in\R\setminus\{0\}\).
The Deligne cup product in the quotient of
Lemma~\ref{lem:geometricdeligne} is represented by
\[
 (it)(2\pi i)h=-2\pi t\,h.
\]
This is nonzero and real, whereas \(H^2(Q_{\C},\R(3))\)
is the imaginary real line through \(h\).
The same calculation can be made on \(\PP^1\): the line bundle
in \eqref{eq:geometricpic} is the pullback of a ruling, and
cup with its Chern class multiplies the point coefficient by
\(2\pi i\). Pullback along the affine-line bundle gives the
displayed calculation on \(Q\).

It follows from regulator naturality and
Lemma~\ref{lem:geometricdeligne} that
\(r_3(\gamma|_{Q_s})\ne0\) for every \(s\in S\).
If an integer \(N>0\) killed \(\gamma|_T\), continuity
\eqref{eq:geometriccontinuity} would kill \(N\gamma\) on some
\(Q_s\). Its regulator there is a nonzero element of a real
vector space, a contradiction. Thus \(\gamma|_T\) has infinite
order.

\section{Flat base change and the conclusion}

The map \(A\to B[1/u]\) is flat: it is a finite free coefficient
extension followed by localization. The divisor \(x=0\) pulls
back to the disjoint union
\[
 (D_+\cap T)\amalg(D_-\cap T),
\]
since \(2u\) is invertible on \(T\).
The original coefficient homomorphism in
\eqref{eq:geometriccoefficient} sends \(u\) to \(z\).
Its two restrictions are therefore \(u\mapsto u\) and
\(u\mapsto-u\). Flat base change and the projection formula give
\begin{align*}
 c|_T
 &= (\beta|_E)[\OO_{D_+\cap T}]
       +(\sigma\beta|_E)[\OO_{D_-\cap T}]\\
 &= (\beta|_E)\bigl([\OO_{D_+\cap T}]
                         -[\OO_{D_-\cap T}]\bigr).
\end{align*}
The union is the principal divisor \(V(x)\).
Its resolution by multiplication by \(x\) gives
\[
 [\OO_{D_+\cap T}]+[\OO_{D_-\cap T}]
   =[\OO_{T/(x)}]=0\quad\text{in }K_0(T).
\]
Consequently
\begin{equation}\label{eq:geometricbasechange}
 c|_T=2(\beta|_E)[\OO_{D_+\cap T}]=2\gamma|_T.
\end{equation}
This class has infinite order, and hence so does \(c\).
Its generic image is zero by its support on \(x=0\), as shown
in Section~1. This proves Theorem~\ref{thm:geometricmain}.

{\small
\begin{thebibliography}{9}
\bibitem{Borel}
A.~Borel, \emph{Stable real cohomology of arithmetic groups},
Ann. Sci. \'Ecole Norm. Sup. (4) \textbf{7} (1974), 235--272;
Proposition~12.2.
\href{https://doi.org/10.24033/asens.1269}{doi:10.24033/asens.1269}.

\bibitem{BunkeTamme}
U.~Bunke and G.~Tamme, \emph{Multiplicative differential algebraic
\(K\)-theory and applications}, Annals of \(K\)-Theory
\textbf{1} (2016), no.~3, 227--258;
Remark~2.6, Theorem~2.31 and Section~6.
\href{https://doi.org/10.2140/akt.2016.1.227}{doi:10.2140/akt.2016.1.227}.

\bibitem{Quillen}
D.~Quillen, \emph{Higher algebraic \(K\)-theory I}, in
\emph{Algebraic \(K\)-Theory I: Higher \(K\)-Theories},
Lecture Notes in Math. \textbf{341}, Springer, 1973, 85--147.
\href{https://doi.org/10.1007/BFb0067053}{doi:10.1007/BFb0067053}.

\bibitem{QuillenFinite}
D.~Quillen, \emph{On the cohomology and \(K\)-theory of the general
linear groups over a finite field}, Ann. of Math. (2)
\textbf{96} (1972), 552--586.
\href{https://doi.org/10.2307/1970825}{doi:10.2307/1970825}.

\bibitem{StacksBlowup}
The Stacks Project Authors, \emph{The Stacks Project},
Section~31.33, \emph{Blowing up}, standard charts and tautological
ideal identity.
\href{https://stacks.math.columbia.edu/tag/01OF}{Tag 01OF}.

\bibitem{Weibel}
C.~A.~Weibel, \emph{The \(K\)-book: An Introduction to Algebraic
\(K\)-theory}, Graduate Studies in Math. \textbf{145}, AMS, 2013;
Chapter~IV, Section~6.4 (filtered-colimit continuity).
\href{https://sites.math.rutgers.edu/~weibel/Kbook/Kbook.IV.pdf}{Author's Chapter IV}.
\end{thebibliography}
}
\end{document}
